\chapter{Discussion}
\label{chap:discussion}


\section{Synthesis of Contributions}

This thesis addresses a fundamental challenge in clinical natural language processing: the scarcity of annotated medical data due to confidentiality constraints, high annotation costs, and limited expert availability.
Unlike general-domain NLP, where crowdsourcing platforms such as Amazon Mechanical Turk enable rapid annotation at scale, clinical data confidentiality prohibits exposing patient documents to untrained external annotators.
The same confidentiality constraints forbid the use of remote large language models for annotation, as transmitting clinical text to external APIs violates data protection regulations.
Furthermore, the limited computational infrastructure of most healthcare institutions makes local deployment of large language models impractical, leaving institutions with few viable annotation options.
We propose methods to automate both the generation of synthetic clinical documents and their annotation, enabling the creation of high-quality training datasets without massive human intervention.

\subsection{A Complete Pipeline for Synthetic Medical Data}

The contributions presented in Chapters~\ref{chap:synthetic} through~\ref{chap:french-application} articulate to form an integrated approach to synthetic data creation.
Chapter~\ref{chap:synthetic} introduces a reward-based generation method that produces synthetic clinical documents guided by UMLS concept keywords extracted from real patient documents.
We term this a \emph{facsimile} approach because the synthetic documents are designed to closely reproduce the structure, style, and clinical reasoning patterns of their real counterparts—much like a facsimile copy preserves the essential characteristics of an original document.
This structural fidelity is not merely aesthetic: models trained on synthetic data that closely resembles local clinical documents transfer more effectively to real documents from the same institution.
Unlike pure generation approaches that generate documents from scratch, our seed-based method anchors synthetic documents in real clinical narratives while extracting only non-identifying medical concepts.

Chapter~\ref{chap:privacy} conducts a rigorous privacy analysis of the generated synthetic documents, revealing that while the keyword-based generation pipeline ensures no direct identification occurs, residual re-identification risks persist through the combination and sequence of medical concepts.
We propose and evaluate post-processing mitigation strategies, demonstrating that targeted keyword perturbation significantly reduces linkage attack success rates while maintaining acceptable data utility for certain downstream tasks.

Chapters~\ref{chap:distillation} and~\ref{chap:french-application} complete the pipeline by addressing the annotation bottleneck.
The hybrid approach combining dictionary-based and LLM-based annotations leverages the complementary strengths of each source, achieving superior performance compared to either method alone.
More broadly, rule-based approaches remain valuable components of clinical NLP pipelines, particularly when high precision is required: hand-crafted rules can enforce domain constraints and filter false positives that statistical models might produce.
The French clinical application (Chapter~\ref{chap:french-application}) validates this approach on real-world data from AP-HP, including semi-synthetic documents created by physicians for educational and research purposes, demonstrating that our methodology transfers effectively across languages and clinical contexts.

Together, these contributions establish an end-to-end process: synthetic document generation (Chapter~\ref{chap:synthetic}) produces privacy-conscious clinical text, privacy mitigation (Chapter~\ref{chap:privacy}) enables safe dissemination for certain attacks, and weak supervision (Chapters~\ref{chap:distillation} and~\ref{chap:french-application}) provides automated annotation.
The result is a framework prototype for creating annotated clinical datasets with minimal human intervention.

\subsection{Privacy Considerations in Practice}

A central theme across this work is the commitment to privacy preservation alongside the honest acknowledgment of its limitations.
The generation pipeline incorporates privacy-by-design principles: only UMLS concept keywords cross the boundary from private to public environments, and the reward mechanism relies solely on numerical scores rather than textual data.
However, the privacy analysis in Chapter~\ref{chap:privacy} reveals that architectural safeguards alone do not eliminate re-identification risks.

The empirical evaluation demonstrates that reward-based training, while improving document quality, inadvertently amplifies re-identification vulnerabilities by teaching the model to reproduce patient-specific patterns.
This finding—that preference optimization based on similarity to real documents reinforces identifying characteristics—represents an important insight for the community developing similarity-based reward mechanisms.

The post-processing keyword perturbation mitigation achieves dramatic reductions in linkage attack success rates, yet discrimination attacks remain highly effective even after mitigation.
This persistent vulnerability underscores a fundamental lesson: there exists no perfect solution that guarantees both complete privacy and full data utility.
This work demonstrates theoretically that significant room for improvement remains in privacy-preserving synthetic data generation.
However, our results also show that under certain constraints—such as deployment in semi-closed systems where synthetic datasets are shared only within controlled research environments rather than published openly—the approach can provide acceptable privacy-utility tradeoffs adapted to the data provider's security requirements.

What we can guarantee is the absence of direct identification: the keywords contain no personally identifiable information such as names, dates, or locations.
What we cannot guarantee is zero re-identification risk, though this risk can be substantially reduced through post-processing.

This raises a practical question: if facsimile documents still carry re-identification risks, can they be submitted to distant LLMs for annotation or other processing?
The answer depends on the data provider's requirements and applicable regulations.
Several deployment configurations can be negotiated with institutions based on their acceptable risk levels: semi-closed access where synthetic datasets are shared within controlled research consortia, post-processing with keyword perturbation to meet specific privacy thresholds, completely open publication when sufficient mitigations have been applied, or transmission to secured infrastructure hosting larger LLMs under appropriate data agreements.
This range of options illustrates the flexibility of synthetic data: institutions can calibrate the privacy-utility tradeoff to their specific constraints rather than being limited to a single dissemination model.

This honest assessment of capabilities and limitations provides a realistic foundation for deployment decisions in healthcare settings.

\subsection{Validation Across Languages and Clinical Contexts}

The effectiveness of our approach is validated not only on public datasets (MIMIC-III) but also on real-world clinical data and across multiple languages.
The weak supervision experiments in Chapter~\ref{chap:distillation} demonstrate strong performance on multilingual clinical entity extraction using the E3C dataset, covering English, Spanish, French, Italian, and Basque.
The French clinical application (Chapter~\ref{chap:french-application}) provides further validation on AP-HP medical reports, demonstrating that the reward-based generation approach transfers effectively to French clinical documentation and ICD-10 classification tasks.

The substantial gap between public datasets (MIMIC-III) and real clinical documents in terms of structure, terminology, and narrative flow motivates the facsimile approach.
The consistent patterns observed across MIMIC-III (English, ICD-9) and AP-HP (French, ICD-10) experiments—despite differences in language, coding systems, and model architectures—suggest that the underlying methodology captures language-agnostic principles of clinical document structure.


\section{Synthetic Document Generation}


\subsection{The Facsimile Approach}

% TODO: Move from §3.1 "Facsimile vs pure generation" (KnowledgeSG comparison)
% TODO: Move from §2.2 "Dependence on seed documents" (scalability + workarounds)


\subsection{Reward-Based Training}

% TODO: Move from §2.1 "Reward-based generation" (privacy-by-design, effectiveness)
% TODO: Move from §2.2 "Reward-based training amplifies privacy"
% TODO: Keep shortened from §2.2 "Discovery of privacy risks came late" (2-3 sentences)
% TODO: Keep shortened from §2.2 "Computational costs" (1 paragraph)


\section{Privacy Preservation}


\subsection{Keyword Perturbation: Strengths and Limits}

% TODO: Move from §2.1 "Keyword perturbation: targeted privacy mitigation"
% TODO: Move from §2.2 "Post-processing mitigation does not eliminate all risks"
% TODO: Move from §3.1 "Empirical privacy mitigation vs formal DP"
% TODO: Move from §3.1 "Rigorous privacy evaluation"


\subsection{Privacy-Utility Tradeoffs in Practice}

% TODO: Move from §4.2 "Conversational and QA applications"
% TODO: Move from §4.2 "Clinical coding and terminology-sensitive"
% TODO: Move from §4.2 "Semi-closed vs open publication"


\section{Weak Supervision for Annotation}


\subsection{Hybrid Dictionary-LLM Approach}

% TODO: Move from §2.1 "Hybrid weak supervision"
% TODO: Condense from §3.2 "Hybrid weak supervision for multilingual" (remove redundancy)
% TODO: Fold from §2.2 "Generalization across languages" (as limitation)


\subsection{Knowledge Distillation}

% TODO: Move from §3.2 "Knowledge distillation in resource-constrained"
% TODO: Move from §3.2 "Positioning as draft annotation"


\section{Research Perspectives}

% TODO: Move from discussion.tex §5 (keep as is)


\subsection{Short-term Extensions}

% TODO: Move from discussion.tex §5.1


\subsection{Long-term Directions}

% TODO: Move from discussion.tex §5.2
